% Source: sections/06_experiments.tex (structure). Numbers: AmphionEval records
% (protocol v4) of ts-robust-064-avg3 and of the two CocktailASR-1 modes, pooled over
% the 248 subsets of Table 1 (v12) and the 220 English subsets (en-v1). The English
% CocktailASR-1 predictions were rescored with the same Whisper normalizer version.
% Cascades (WeSep or LauraTSE + Qwen3-ASR-1.7B): full AmphionEval external-worker runs on
% both sets (2026-11-22, same protocol and normalizer), pooled with the same script.
\section{Experiments}
\label{sec:experiments}
\subsection{Protocol and baselines}
We use one set of \model weights across all subsets in both languages, with the first 5\,s of enrollment and the same normalization for all systems.
\model uses greedy decoding with at most 512 output tokens.
The end-to-end (E2E) TS-ASR baseline is CocktailASR-1~\cite{zhang2026xiaomi} in its standard and chain-of-thought modes, denoted CocktailASR-1 and CocktailASR-1-CoT, using the publicly available \texttt{Ease3/Xiaomi-CocktailASR-1} weights and their accompanying inference code with bf16 and greedy decoding, at most 256 output tokens in Mandarin and 512 in English.
Its report evaluates enrollment-conditioned transcription and rejection of absent targets in both Chinese and English.
Baseline scores are measured on the same audio and scoring protocol, not copied from their reports.

We also evaluate two extract-then-recognize cascades.
A target speech extraction (TSE) model receives the same enrollment and the mixture, and Qwen3-ASR-1.7B transcribes its output with greedy decoding.
WeSep~\cite{wang2024wesep} uses its public English BSRNN model with an ECAPA speaker encoder.
LauraTSE~\cite{tang2025lauratse}, a generative extractor built on a decoder-only language model, uses offline inference with top-25 sampling and a fixed seed.
We denote them WeSep + Qwen3-ASR-1.7B and LauraTSE + Qwen3-ASR-1.7B.
Neither cascade has a separate rejection step: an absent target counts as rejected only when the transcript is empty.

The CocktailASR-1 report uses reference crops of 1--4\,s.
On an earlier benchmark version, a fixed 3\,s crop and that random-crop protocol produced similar target-absent FA, although group CER differed by up to about seven points.
Our main comparison holds enrollment length fixed across systems.

\subsection{One model across scenarios}
Table~\ref{tab:overall} reports all three metrics.
In Mandarin, \model achieves 13.1\% CER against 33.4\% for CocktailASR-1, with an ER 0.6 points higher (3.3\% versus 2.7\%).
Thus the lower FA is not explained by a large increase in completely suppressed target utterances.
Overall FA is 2.8\% versus 67.8\%; CocktailASR-1-CoT reduces FA to 61.9\% but increases CER to 39.6\%.
The cascades rarely return empty transcripts.
WeSep + Qwen3-ASR-1.7B has the lowest ER, 0.3\%, but 87.0\% FA and 62.0\% CER; LauraTSE + Qwen3-ASR-1.7B has 64.9\% FA and 80.5\% CER.

On real conversations with clean enrollments, CER is 9.2\% versus 30.8\%.
For individual real corpora with clean enrollments, \model has FA between 1\% and 7\%, compared with 74--100\% for CocktailASR-1.
For example, AliMeeting far-field CER is 6.3\% versus 38.2\%, and FA is 1.0\% versus 100.0\%.
When the enrollment contains another voice, our CER rises from 10.4\% to 15.0\% on AliMeeting, AISHELL-4, and web-conversation windows, and from 7.3\% to 18.6\% on the paired AliMeeting windows; with such enrollments, FA is 0.7\%.
These are whole-system comparisons: the study does not separately establish the contributions of architecture, training data, or parameter count.

\subsection{What restricted evaluations miss}
Table~\ref{tab:diagnostics} contrasts common evaluation conditions with cases encountered during continuous listening.
Single-talker real recordings barely distinguish the systems.
Adding silence-only negatives gives CocktailASR-1 a 10.8\% FA, while negatives containing other speakers reveal much higher FA.
With two to four interferers and the target present, its CER reaches 65.7\%, versus 18.5\% for \model.
Both target presence and conversation structure therefore affect the comparison.
The cascades follow the same pattern more strongly.
On single-talker real recordings, WeSep + Qwen3-ASR-1.7B has the lowest CER, 5.4\%, but with two to four interferers both cascades exceed 190\% CER, and with only others speaking their FA is 70--99\%.

Paired enrollment tests expose a remaining weakness.
For 6--18\,s absent-target segments, replacing a random enrollment with a similar-sounding person's enrollment raises our FA from 0.0\% to 10.8\%.
For a single short phrase, it rises from 3.0\% to 16.7\%.
Only 1\,s of enrollment speech also raises FA to 18.0\%.
Short, similar voices remain difficult despite the low overall FA.

Real near/far-field pairs qualify synthetic robustness conclusions.
With near-field enrollment, switching the mixture to far field increases our CER from 9.4\% to 18.9\% synthetically, but from 4.6\% to 7.6\% on the paired AliMeeting windows.
Synthetic subsets identify sensitivity; real windows measure its magnitude under actual recording conditions.

\subsection{English counterpart}
The English columns of Table~\ref{tab:overall} apply the same protocol to the English counterpart.
\model has 13.6\% WER against 26.3\% for CocktailASR-1, and an ER of 3.7\% against 5.0\%; CocktailASR-1-CoT has the lowest ER, 0.6\%.
FA is 6.3\% versus 50.6\%, and 81.8\% and 58.7\% for WeSep + Qwen3-ASR-1.7B and LauraTSE + Qwen3-ASR-1.7B.
With similar-sounding absent speakers, our FA is 20.2\%, against 13.9\% in Mandarin.

Real English conversations are harder than Mandarin ones: FA is 19.5\%, against 2.4\% in Mandarin.
With clean enrollments, our FA ranges from 2.4\% on ICSI and 5.2\% on NOTSOFAR-1 to 28.6\% on DiPCo and 69.3\% on CHiME-6, while CocktailASR-1 has 92.7--100\% on every corpus.

\subsection{Limitations}
Both language versions are mostly synthetic, and the English real section has only 1{,}301 samples.
Real data cover meetings, calls, and web conversations; web references contain pseudo-label errors.
Training uses the training splits of the Mandarin real meeting corpora and also includes recordings from all five English real corpora, so the real results measure in-domain behavior.
The same team designed training scenes and benchmark conditions.
AISHELL-1 test text overlaps with AISHELL-2 training text, and MagicData-RAMC speaker-ID inconsistencies may allow speaker overlap.
The benchmark therefore does not establish unseen-domain or fully speaker-disjoint generalization.
There is no controlled architectural ablation or formal significance test.
The cascades use public extraction checkpoints trained on LibriSpeech-based mixtures (LauraTSE) and VoxCeleb (WeSep)~\cite{tang2025lauratse}, without adaptation to Mandarin or to our training data; they compare available systems, not extraction as an approach.
An empty transcript is not an identity guarantee, and recordings longer than the request limit require segmentation.
